\begin{pastexam}[label={px:nchu-111-is-10}]{中興資工所乙組 111 學年度 第 21 題}
（問答）

{[3\%]} Please explain what is the NP (Nondeterministic Polynomial) problem.
\end{pastexam}

\begin{pastsoln}
NP 是一類\textbf{判定問題}（答案只有是或否）的集合（資結書第~28.2~節）：一個判定問題屬於 NP，若答案是
「是」的每個輸入，都存在一個長度是輸入大小的多項式的\textbf{證據}（certificate），而且有一個多項式時間的
驗證演算法，拿到輸入和證據就能檢查答案確實是「是」。 簡單說：答案是「是」時，可以在多項式時間內驗證。

名稱裡的「nondeterministic」來自等價的定義：非決定性的機器可以一次「猜」出證據，再用多項式時間驗證，
所以 NP 是非決定性機器在多項式時間內能解的問題。 它\textbf{不是}「non-polynomial」（不能在多項式時間內
解）的意思。

例子：
\begin{itemize}
\item 漢米爾頓迴路：給一個圖，問有沒有經過每個點恰一次的環。 證據是那個環（點的排列），驗證時檢查每個點
恰出現一次、相鄰的點之間有邊，$O(n^2)$ 以內。
\item SAT：給一個布林公式，問能不能讓它為真。 證據是一組真假指派，代進去算一次就能驗證。
\end{itemize}
P（能在多項式時間內\textbf{解}的問題）都在 NP 裡，因為直接解一次就算是驗證，不需要證據（資結書定理~28.1）。
NP 裡最難的是 NP 完備問題：每個 NP 問題都能在多項式時間內歸約到它（資結書第~28.4~節）。 P 是否等於 NP
至今沒有人知道。
\end{pastsoln}
